\documentclass{article}
\usepackage[utf8]{inputenc}
\usepackage{amsmath}
\usepackage{amssymb}
\usepackage{graphicx}
\usepackage{tabularx}
\usepackage{booktabs}
\usepackage{caption}
\usepackage[english]{babel}
\usepackage{setspace}
\usepackage{geometry}
\usepackage{float}
\usepackage{url}
\usepackage[justification=centering]{caption}

% APA style formatting
\geometry{a4paper, margin=1in}
\setstretch{1.5}
\captionsetup[table]{labelfont=bf, justification=raggedright, singlelinecheck=false}
\captionsetup[figure]{font=small, labelfont=bf, justification=raggedright, singlelinecheck=false}

\title{ Machine Learning for Customer Churn Prediction}
\author{Nelson David Posso Suarez \\ Javier Alejandro Penagos Hernandez \\ Nicolas Castro Rivera }
\date{}

\begin{document}

\maketitle

\section{Summary}

This work formulates the customer churn retention problem as a Markov Decision Process (MDP).
The objective is to train a Reinforcement Learning (RL) agent to select the best retention strategy
for a customer with a probability of churn. Although the available dataset does not contain labeled
data describing the outcome of different retention strategies, it provides sufficient customer and
service information to establish an environment where the agent can learn which actions can reduce
the probability of churn.



\section{MDP Formulation}
The MDP is composed of a state space containing customer features, an action set representing
different retention strategies, stochastic transitions based on the customer's acceptance of an
offer, and a reward function based on the reduction of the predicted churn probability and the
cost of the selected action. A Deep Q-Network (DQN) algorithm is used as the RL agent to determine
the action with the highest expected reward.

\subsection{Problem Statement}

The objective of the MDP is to determine the best retention strategy for a customer according to
their characteristics and current probability of churn. The RL agent receives the customer's state
and selects an action from a set of possible retention strategies. The selected action can be
accepted or rejected by the customer, introducing stochasticity into the environment.

The customer's churn probability is obtained using a supervised learning model. This probability
is used by the RL environment to evaluate whether the selected retention strategy improves the
customer's situation. The objective of the agent is therefore to reduce the probability of churn
while considering the cost associated with some retention strategies.

\subsection{State Space}

The state consists of the features that describe the customer and their current service configuration.
The original Telco Customer Churn dataset contains 19 customer and service features, without count the customer id. In addition,
2 features are created through feature engineering, and 1 additional feature is generated by the
supervised learning model. This feature represents the predicted probability of the customer
being churn. 

For this reason, each state is represented as a vector of 22 elements:

\begin{center}
\texttt{State = [19 original features, 2 engineered features, churn probability]}
\end{center}


\subsection{Action Set}

The action set contains 15 possible retention strategies. Each action modifies specific features
of the customer's state. The actions are divided into three categories: contract plan changes,
monthly charge discounts, and new service offers.

\subsubsection{Change the Contract Plan}

Customers with a month-to-month contract usually have a higher probability of churn than customers
with one-year or two-year contracts. Therefore, the agent can offer a change in the customer's
contract duration. These actions have zero cost.

\begin{enumerate}
\item Offer change to a One Year contract.
\item Offer change to a Two Years contract.
\end{enumerate}

\subsubsection{Discount in the Monthly Charge}

High monthly charges can increase the probability of customer churn. For this reason, the agent
can offer a discount of 5\%, 10\%, or 15\% on the customer's monthly charge.

These actions generate a cost for the company. The cost is proportional to the size of the
discount and is defined as the discount percentage multiplied by one hundred.

\begin{enumerate}
\setcounter{enumi}{2}
\item Offer a 5\% discount on the Monthly Charges.
\item Offer a 10\% discount on the Monthly Charges.
\item Offer a 15\% discount on the Monthly Charges.
\end{enumerate}

\subsubsection{Offer a New Service}

The company provides 10 services that can be offered to customers. When a service is accepted,
it is added to the customer's plan and its estimated price is added to the Monthly Charges.
The service prices were estimated from the dataset to obtain approximate values in USD.

There are two restrictions for these actions. Multiple Lines cannot be offered if the customer
does not have Phone Service. The other services cannot be offered if the customer does not have
Internet Service. These actions have zero cost.

\begin{enumerate}
\setcounter{enumi}{5}
\item Offer Phone Service (USD 20).
\item Offer Multiple Lines (USD 5).
\item Offer Internet Service DSL (USD 25).
\item Offer Internet Service Fiber Optic (USD 50).
\item Offer Online Security (USD 5).
\item Offer Online Backup (USD 5).
\item Offer Device Protection (USD 5).
\item Offer Tech Support (USD 5).
\item Offer Streaming TV (USD 10).
\item Offer Streaming Movies (USD 10).
\end{enumerate}

\subsection{Transition Dynamics}

The selected action directly affects the customer's state. However, before modifying the state,
the environment determines whether the customer accepts or rejects the offered strategy. This
makes the transition stochastic because the same action can produce different outcomes.

The acceptance probabilities are defined according to the type of retention strategy:

\begin{itemize}
\item Change the contract plan: 85\% acceptance and 15\% rejection.
\item Discount in the monthly charge: 95\% acceptance and 5\% rejection.
\item Offer a new service: 90\% acceptance and 10\% rejection.
\end{itemize}

The transition function can be represented as:

\begin{equation}
s_{t+1} = f(s_t, a_t)
\end{equation}

The acceptance or rejection of an action is stochastic. However, once an action is accepted,
the modification applied to the customer's state is deterministic. Therefore, the same state
and the same accepted action produce the same resulting state.

\subsection{Reward Function}

The reward function measures the benefit obtained from applying a retention strategy. It is based
on the difference between the customer's churn probability before and after the action, considering
the cost associated with the selected strategy.

The reward is calculated using the following function:

\begin{equation}
Reward = P_{before} - P_{after} - Cost
\end{equation}

The resulting values are multiplied by one hundred.

\begin{itemize}

\item \texttt{P\_before}: Churn probability predicted by the supervised model at the beginning
of the episode.

\item \texttt{P\_after}: Churn probability predicted by the supervised model after the action
is accepted and the customer's state is modified.

\item \texttt{Cost}: Cost associated with the selected retention strategy.

\item \texttt{Reward}: Benefit obtained from reducing the customer's probability of churn while
considering the cost of the action.

\end{itemize}

A reduction in the predicted churn probability produces a positive reward, while an increase in
the probability or a high action cost reduces the resulting reward.

\subsection{Episode }

An episode starts with a customer from the Telco Customer Churn dataset. The customer's initial
state is provided to the supervised learning model, which predicts the probability of churn.
This information is then provided to the RL agent.

The agent uses the DQN algorithm to select the retention strategy with the highest expected reward.
The selected action is then evaluated by the environment, which determines whether the customer
accepts or rejects the offer according to the corresponding acceptance probability.

If the customer accepts the offer, the action modifies the customer's state. The supervised model
then predicts the new probability of churn, which is compared with the probability before the action
to calculate the reward. The episode terminates after this reward is obtained.

If the customer rejects the offer, the customer's state remains unchanged and the reward is zero.
The episode then terminates.


\section{Environment}

\section{Agent algorithm}

\section{Hyperparameter}

\section{Sample trajectory visualizations}

\section{Discussion}

\section{Conclusion}


\begin{thebibliography}{9}

\bibitem{ahmad2019}
Ahmad, A. K., Jafar, A., \& Aljoumaa, K. (2019). Customer churn prediction in telecom using machine learning in big data platform. \textit{Journal of Big Data}, 6, 28.

\bibitem{verbeke2012}
Verbeke, W., Dejaeger, K., Martens, D., Hur, J., \& Baesens, B. (2012). New insights into churn prediction in the telecommunication sector: A profit driven data mining approach. \textit{European Journal of Operational Research}, 218(1), 211--229.

\bibitem{breiman2001randomforests}
Breiman, L. (2001). Random forests. \textit{Machine Learning}, 45(1), 5--32.

\bibitem{chen2016xgboost}
Chen, T., \& Guestrin, C. (2016). XGBoost: A scalable tree boosting system. In \textit{Proceedings of the 22nd ACM SIGKDD International Conference on Knowledge Discovery and Data Mining} (pp. 785--794).

\bibitem{akiba2019optuna}
Akiba, T., Sano, S., Yanase, T., Ohta, T., \& Koyama, M. (2019). Optuna: A next-generation hyperparameter optimization framework. In \textit{Proceedings of the 25th ACM SIGKDD International Conference on Knowledge Discovery \& Data Mining} (pp. 2623--2631).

\bibitem{paszke2019pytorch}
Paszke, A., Gross, S., Massa, F., Lerer, A., Bradbury, J., Chanan, G., Killeen, T., Lin, Z., Gimelshein, N., Antiga, L., Desmaison, A., Kopf, A., Yang, E., DeVito, Z., Raison, M., Tejani, A., Chilamkurthy, S., Steiner, B., Fang, L., Bai, J., \& Chintala, S. (2019). PyTorch: An imperative style, high-performance deep learning library. \textit{Advances in Neural Information Processing Systems}, 32.

\bibitem{chawla2002smote}
Chawla, N. V., Bowyer, K. W., Hall, L. O., \& Kegelmeyer, W. P. (2002). SMOTE: Synthetic minority over-sampling technique. \textit{Journal of Artificial Intelligence Research}, 16, 321--357.

\bibitem{davis2006precision}
Davis, J., \& Goadrich, M. (2006). The relationship between precision-recall and ROC curves. In \textit{Proceedings of the 23rd International Conference on Machine Learning} (pp. 233--240).

\end{thebibliography}

\end{document}
